\section*{Chapter \ref{ch:g10:sensors-actuators-control} --- Sensors, Actuators and Control Programs}

\begin{solution}{exo:g10:sensors-actuators-control:1}
$-20 + 300 \times 80/1023 \approx 3.46$~°C and $-20 + 700 \times 80/1023 \approx 34.74$~°C. $0$~°C would be
$20 \times 1023/80 \approx 255.8$: the reading $256$ ($0.02$~°C) is closest, nearer than $255$ ($-0.06$~°C).
\end{solution}

\begin{solution}{exo:g10:sensors-actuators-control:2}
\textit{amber} stays \textit{amber}; \textit{red} becomes \textit{green}. The fastest way to red is \textit{button},
\textit{tick}, \textit{tick}, \textit{tick}: green, green wait, amber, red.
\end{solution}

\begin{solution}{exo:g10:sensors-actuators-control:3}
Three windows of five: $100.4/5 = 20.08$, $99.9/5 = 19.98$ and $100.2/5 = 20.04$. The readings spread over
$0.8$ degree, the averages over $0.1$.
\end{solution}

\begin{solution}{exo:g10:sensors-actuators-control:4}
$6 \times 4 = 24$ entries, one per pair (\omterm{def:g10:sensors-actuators-control:state-machine}{state}, \omterm{def:g10:sensors-actuators-control:state-machine}{event}). The traffic light has $4$ \omterm{def:g10:sensors-actuators-control:state-machine}{states} and $2$ \omterm{def:g10:sensors-actuators-control:state-machine}{events}: its table has
$8$ entries, all present, so it is complete.
\end{solution}

\begin{solution}{exo:g10:sensors-actuators-control:5}
\code{return -40 + raw * 125 / 1023} (the range is $85 - (-40) = 125$ degrees). The step is $125/1023 \approx 0.122$
degree. \code{exercises/controllers.py} has the same function with the range as parameters.
\end{solution}

\begin{solution}{exo:g10:sensors-actuators-control:6}
\code{TABLE[("amber", "button")]} raises a \code{KeyError} and the program stops: the controller is dead, and the
light stays as it was, or goes dark, while cars and pedestrians rely on it. The function \code{next_state} of the
exercise file returns \code{table.get((state, event), state)}, so that a missing pair means ``no change''. Better still, a test checks that the table has all $8$ pairs.
\end{solution}

\begin{solution}{exo:g10:sensors-actuators-control:7}
\code{if t < 18: return True}, \code{if t > 21: return False}, else \code{return heater}: the function receives the
current \omterm{def:g10:sensors-actuators-control:state-machine}{state} of the heater and returns the new one (\code{exercises/controllers.py}).
\end{solution}

\begin{solution}{exo:g10:sensors-actuators-control:8}
The plain thermostat switches at every crossing of $19.5$: on at $19.4$, off at $19.6$, on at $19.3$, off at
$19.7$, on at $17.8$, off at $19.8$, on at $19.4$, off at $19.6$: $8$ times. The greenhouse thermostat switches on at
$17.8$ and off at $21.2$: $2$ times. The tests check both counts.
\end{solution}

\begin{solution}{exo:g10:sensors-actuators-control:9}
green, green wait, amber, amber (the button changes nothing), red, green, green wait, amber. The cars wait one tick at
red, the time given to the pedestrians.
\end{solution}

\begin{solution}{exo:g10:sensors-actuators-control:10}
\textit{call1} (or \textit{call2}): the new version gives $(0, \text{open}, \text{up}, 1)$, and \code{safe} returns
\code{False}: the cabin would leave with its door open.
\end{solution}

\begin{solution}{exo:g10:sensors-actuators-control:11}
Two nested loops over the \omterm{def:g10:sensors-actuators-control:state-machine}{states} and the \omterm{def:g10:sensors-actuators-control:state-machine}{events}, keeping the pairs with \code{safe(s)} true and
\code{safe(step(s, e))} false (\code{exercises/controllers.py}). On the lift's $12$ reachable \omterm{def:g10:sensors-actuators-control:state-machine}{states} the list is
empty. With the new version it has $6$ pairs: an open door at any floor, with a call to either other floor.
\end{solution}

\begin{solution}{exo:g10:sensors-actuators-control:12}
The controller returns \code{True} only in its first branch, which needs a reading below $a$, and \code{False} only in
its second, which needs a reading above $b$; otherwise it returns the current \omterm{def:g10:sensors-actuators-control:state-machine}{state}. So an off-to-on switch happens
only at a reading below $a$ and an on-to-off switch only at a reading above $b$. Between two switchings on there is a
switching off, so some reading in between was above $b$, and the next switching on needs a reading below $a$ again:
the readings crossed the whole band, from below $a$ to above $b$ and back.
\end{solution}

\begin{solution}{exo:g10:sensors-actuators-control:13}
The \omterm{def:g10:sensors-actuators-control:state-machine}{state} is (cars inside, barrier position). On \textit{car at entry}, if there are fewer than $3$ cars, the count
goes up by one; on \textit{car at exit}, if there is a car, it goes down by one (\code{exercises/controllers.py}).
Four arrivals give $1$, $2$, $3$, $3$ (the fourth car is refused), and one departure gives $2$.
\end{solution}

\begin{solution}{exo:g10:sensors-actuators-control:14}
From $(0, \text{closed}, \text{stop}, \text{None})$: $(0, \text{open}, \text{stop}, \text{None})$ by \textit{open}
or \textit{call0}; $(0, \text{closed}, \text{up}, 1)$ by \textit{call1}; then $(1, \text{open}, \text{stop},
\text{None})$ by \textit{arrive}; $(1, \text{closed}, \text{stop}, \text{None})$ by \textit{close}; and
$(1, \text{closed}, \text{down}, 0)$ by \textit{call0}, whose \textit{arrive} leads back to floor $0$ with the door
open. Six \omterm{def:g10:sensors-actuators-control:state-machine}{states}; the open ones all have the motor stopped, so all are safe.
\end{solution}

\begin{solution}{exo:g10:sensors-actuators-control:15}
\code{input()} waits until a line is typed: while nobody types, the loop is stuck, the thermometer is not read and the
heater stays as it was, perhaps on for hours. The loop must never wait for a person: at each turn it checks whether
a command has arrived (a button is a \omterm{def:g10:sensors-actuators-control:sensor}{sensor} too, read like the others), handles it if there is one, and goes on with
its readings either way.
\end{solution}

\begin{solution}{pb:g10:sensors-actuators-control:1}
\textbf{1.} Floor $3$, door $2$, motor $3$, target $4$ ($0$, $1$, $2$ or \code{None}): at most
$3 \times 2 \times 3 \times 4 = 72$ tuples.

\textbf{2.} \textit{call0}, \textit{call1}, \textit{call2}, \textit{open}, \textit{close}, \textit{arrive}.

\textbf{3.} The target is the floor the cabin is going to; when it is stopped it is going nowhere, and the target is
\code{None}.

\textbf{4.} Those with the door open and the motor going up or down: $3 \times 2 \times 4 = 24$ of them.

\textbf{5.} \textit{call0} and \textit{open}: $(0, \text{open}, \text{stop}, \text{None})$; \textit{call1}:
$(0, \text{closed}, \text{up}, 1)$; \textit{call2}: $(0, \text{closed}, \text{up}, 2)$; \textit{close} and
\textit{arrive}: no change.

\textbf{6.} \textit{arrive} gives $(2, \text{open}, \text{stop}, \text{None})$; the five others change nothing, since
the motor is running.

\textbf{7.} \textit{open} (only with the motor stopped), a call to the current floor (only with the motor stopped),
and \textit{arrive} at the target, which stops the motor in the same \omterm{def:g10:sensors-actuators-control:state-machine}{transition}. Each leaves the motor stopped.

\textbf{8.} Only a call to another floor, which requires the door closed and the motor stopped.

\textbf{9.} $12$.

\textbf{10.} Stopped: each floor with the door open or closed, $6$ \omterm{def:g10:sensors-actuators-control:state-machine}{states}. Moving: up from $0$ to $1$, up from $0$ to
$2$, up from $1$ to $2$, down from $2$ to $1$, down from $2$ to $0$, down from $1$ to $0$: $6$ \omterm{def:g10:sensors-actuators-control:state-machine}{states}.

\textbf{11.} A call to the floor where the cabin is opens the door instead of starting the motor, and \textit{arrive}
at the target stops the motor: the cabin never moves towards the floor it is on.

\textbf{12.} The set found contains the start, and every \omterm{def:g10:sensors-actuators-control:state-machine}{event} applied to any of its \omterm{def:g10:sensors-actuators-control:state-machine}{states} gives a \omterm{def:g10:sensors-actuators-control:state-machine}{state} of the set:
whatever the sequence of \omterm{def:g10:sensors-actuators-control:state-machine}{events}, the lift moves from \omterm{def:g10:sensors-actuators-control:state-machine}{state} to \omterm{def:g10:sensors-actuators-control:state-machine}{state} inside the set, and every \omterm{def:g10:sensors-actuators-control:state-machine}{state} of the set is
safe.

\textbf{13.} Yes: $(0, \text{open}, \text{stop}, \text{None})$ is in the set, so everything reachable from it is in
the set too.

\textbf{14.} Call $2$ from floor $0$, door closed: the call and two \textit{arrive}, $3$ \omterm{def:g10:sensors-actuators-control:state-machine}{events}. Call $0$ from floor
$2$, door open: \textit{close}, the call and two \textit{arrive}, $4$ \omterm{def:g10:sensors-actuators-control:state-machine}{events}.

\textbf{15.} $(3 + 4 + 3 + 3 + 3)/5 = 3.2$.

\textbf{16.} $2, 0, 2, 0, 2$ needs $3, 4, 4, 4, 4$ \omterm{def:g10:sensors-actuators-control:state-machine}{events}, a mean of $3.8$. $1, 1, 1, 1, 1$ needs $2$ each, a mean
of $2$.

\textbf{17.} The target: it would become a set of floors still to serve, and \textit{arrive} would stop at any floor
in that set.

\textbf{18.} The number of \omterm{def:g10:sensors-actuators-control:state-machine}{states} grows with the floors and the targets, to a few hundred for one target, and much
more with a set of calls; a program still lists them all in a fraction of a second, which is why such checks are
done by computer.

\textbf{19.} The lift has $12$ reachable \omterm{def:g10:sensors-actuators-control:state-machine}{states} among $72$ possible tuples, and none has the door open with the motor
running: the door opens only in \omterm{def:g10:sensors-actuators-control:state-machine}{transitions} that leave the motor stopped, and the motor starts only from a \omterm{def:g10:sensors-actuators-control:state-machine}{state}
with the door closed. Since every \omterm{def:g10:sensors-actuators-control:state-machine}{event} leads from a reachable \omterm{def:g10:sensors-actuators-control:state-machine}{state} to a reachable \omterm{def:g10:sensors-actuators-control:state-machine}{state}, this check of all $12$
proves that the door is never open while the cabin moves.
\end{solution}
